\section{Introducción y arquitectura general}

\subsection{Objetivos y enfoque de la solución}

El proyecto ``Habla con tu dinero'' es nuestra solución al reto propuesto por 
la Cátedra de Inteligencia Artificial de la Universidad de Granada junto a Unicaja. 

Con él, buscamos que el usuario pueda realizar consultas y operaciones financieras de 
manera sencilla e interactiva.

Para garantizar un funcionamiento fiable, se ha separado la interpretación de
las consultas de la ejecución de las operaciones. Mientras que el modelo de lenguaje 
facilita la interacción con el usuario, el \emph{backend}
se encarga de obtener los datos, realizar los cálculos y ejecutar las
operaciones correspondientes.

\subsection{Arquitectura técnica}

La aplicación se divide en tres bloques principales: la interfaz de usuario,
el \emph{backend} y la base de datos.



\begin{itemize}

    \item \textbf{Interfaz de usuario (Frontend):} Se ha desarrollado utilizando
    HTML y JavaScript. La comunicación con
    el servidor se realiza mediante WebSockets, manteniendo una conexión abierta
    durante la interacción. Esto permite recibir las respuestas de forma
    progresiva y actualizar la interfaz dinámicamente. Además, se pueden mostrar
    elementos como gráficos, tablas y mensajes dentro de la propia conversación.

    \item \textbf{Motor principal (Backend):} Se ha desarrollado en Python
    utilizando FastAPI. Se encarga de recibir las peticiones del usuario, gestionar 
    el contexto de la conversación y coordinar los diferentes componentes del sistema. 
    También permite resolver directamente determinadas consultas sin utilizar el modelo de
    lenguaje. Por ejemplo, una consulta sobre el saldo de una cuenta puede
    resolverse directamente desde el backend, reduciendo el tiempo de respuesta.

    \item \textbf{Base de datos y seguridad:} La información de la aplicación
    se almacena en una base de datos SQLite. El acceso a los datos se gestiona
    desde el backend, que determina qué operaciones puede realizar cada parte
    del sistema. Las consultas que únicamente necesitan información utilizan
    acceso de lectura, mientras que las operaciones que modifican datos o
    implican movimientos de dinero son ejecutadas directamente por la lógica de
    la aplicación.

\end{itemize}

Como se muestra en la Figura~\ref{fig:arquitectura}, el 
\emph{backend} se estructura de forma modular. Un controlador 
principal (\texttt{main.py}) gestiona las conexiones y delega el 
flujo en un agente (\texttt{agent.py}). Este agente utiliza un 
catálogo de herramientas (\texttt{tools.py}) para derivar la 
ejecución hacia el motor visual (\texttt{graficos.py}), la 
lógica analítica (\texttt{analitica.py}), las APIs bancarias 
(\texttt{banking\_api.py}) o la base de datos 
(\texttt{database.py}). 

Asimismo, el modelo de lenguaje (por defecto, Qwen3) se ejecuta 
de forma local y privada a través de Ollama, comunicándose 
exclusivamente con el \emph{backend} sin acceso directo a los 
datos.

\begin{figure}[htbp]
\centering
\begin{tikzpicture}[
    font=\small,
    comp/.style={draw=#1, fill=#1!7, rounded corners=2pt, align=center,
                 minimum width=3.2cm, minimum height=1.05cm,
                 inner sep=3pt},
    comp/.default=subsectionblue,
    grupo/.style={draw=gray!55, dashed, rounded corners=5pt, inner sep=7pt},
    titulo/.style={font=\footnotesize\bfseries, text=gray!65!black},
    flecha/.style={-{Stealth[length=2mm]}, thick, draw=gray!70!black},
    doble/.style={{Stealth[length=2mm]}-{Stealth[length=2mm]}, thick,
                  draw=gray!70!black},
    nota/.style={font=\scriptsize, text=gray!55!black, align=center},
]

% --- Navegador -------------------------------------------------------------
\node[comp] (voz)  at (-5.4, 0) {Voz\\[-1pt]\scriptsize STT y TTS};
\node[comp] (chat) at (-1.8, 0) {Chat\\[-1pt]\scriptsize texto por frases};
\node[comp] (vis)  at ( 1.8, 0) {Visuales\\[-1pt]\scriptsize Vega-Lite y tablas};
\node[comp] (pin)  at ( 5.4, 0) {Teclado del PIN\\[-1pt]\scriptsize evento aparte};

\begin{scope}[on background layer]
  \node[grupo, fit=(voz)(pin),
        label={[titulo, anchor=south west]north west:Navegador · \texttt{index.html}}]
        (nav) {};
\end{scope}

% --- Backend ---------------------------------------------------------------
\node[comp] (main) at (-5.4, -3.1) {\texttt{main.py}\\[-1pt]\scriptsize FastAPI};
\node[comp] (ag)   at (-1.8, -3.1) {\texttt{agent.py}\\[-1pt]\scriptsize bucle agéntico};
\node[comp] (tl)   at (-1.8, -4.8) {\texttt{tools.py}\\[-1pt]\scriptsize 6 herramientas};
\node[comp] (gr)   at ( 1.8, -4.8) {\texttt{graficos.py}\\[-1pt]\scriptsize motor visual};

\node[comp] (db)   at (-5.4, -6.5) {\texttt{database.py}\\[-1pt]\scriptsize SQL del LLM};
\node[comp] (api)  at (-1.8, -6.5) {\texttt{banking\_api.py}\\[-1pt]\scriptsize saldo y Bizum};
\node[comp] (an)   at ( 1.8, -6.5) {\texttt{analitica.py}\\[-1pt]\scriptsize recurrencia y previsión};

\begin{scope}[on background layer]
  \node[grupo, fit=(main)(gr)(db)(an),
        label={[titulo, anchor=south east]north east:Backend · Python}]
        (back) {};
\end{scope}

% --- Modelo y datos --------------------------------------------------------
\node[comp=ugrred, minimum height=1.3cm] (llm) at (6.1, -3.1)
      {Ollama\\[-1pt]\scriptsize \texttt{qwen3:8b} en GPU local};

\node[comp=gray, minimum width=10.4cm] (bd) at (-1.8, -8.6)
      {SQLite · \texttt{banco.db}\\[-1pt]\scriptsize 24 meses de movimientos ficticios};

% --- Conexiones ------------------------------------------------------------
\draw[doble] (main.north |- nav.south) -- node[nota, right=2pt] {WebSocket\\eventos JSON} (main.north);

\draw[doble]  (main) -- (ag);
\draw[flecha] (ag) -- (tl);
\draw[flecha] ([yshift=-0.3cm]ag.east) -| node[nota, right, pos=0.75] {tras la\\respuesta} (gr.north);

\draw[doble]  ([yshift=0.2cm]ag.east) -- node[nota, above] {llamada principal (streaming)} ([yshift=0.2cm]llm.west);
\draw[flecha] (gr.east) -| node[nota, below right, pos=0.3] {llamada dedicada} (llm.south);

\draw[flecha] (tl.south) -- ++(0,-0.4) -| (db.north);
\draw[flecha] (tl) -- (api);
\draw[flecha] (tl.south) -- ++(0,-0.4) -| (an.north);

\draw[flecha] (db)  -- node[nota, left]  {solo lectura}               (db  |- bd.north);
\draw[doble]  (api) -- node[nota, right] {\texttt{BEGIN IMMEDIATE}}  (api |- bd.north);
\draw[flecha] (an)  -- node[nota, right] {lectura}                    (an  |- bd.north);

\end{tikzpicture}
\caption{Arquitectura del sistema. El modelo se ejecuta en local y el backend
es el único que accede a la base de datos: el SQL generado por el LLM entra por
una conexión de solo lectura, y las escrituras solo se hacen desde las APIs
bancarias, dentro de una transacción.}
\label{fig:arquitectura}
\end{figure}
